\documentclass[aspectratio=169]{beamer}
\input{header}
\coursecode{JKSSB}

\title{Humanware \& IT Governance}
\subtitle{Lesson 47 --- The Human Element and IT in Governance}
\author{JKSSB Computer Awareness}
\date{\today}

\begin{document}
\frame{\titlepage}

\begin{frame}{What Is Humanware?}
  \begin{itemize}
    \item A computer system needs people as much as it needs machines.
    \item \key{Humanware} (also called \key{liveware} or \key{peopleware}) is
      the human element of a computer system.
    \item It means the people who design, build, operate, manage, and use
      computers.
    \item Humanware is \key{neither hardware nor software}; it is the people
      component.
    \item \muted{Every other part of the system exists to serve the people
      who use it.}
  \end{itemize}
  \begin{block}{Definition}
    Humanware is the people involved with a computer system: its users,
    operators, administrators, programmers, and support staff.
  \end{block}
\end{frame}

\begin{frame}{The People in a Computer System}
  \begin{center}
    \begin{tabular}{p{0.30\textwidth}p{0.58\textwidth}}
      \toprule
      \textbf{Role} & \textbf{What the person does} \\
      \midrule
      User / end user & Uses the computer to do day-to-day work \\
      Operator & Runs the machine and routine jobs \\
      Administrator & Manages systems, networks, databases, security \\
      Programmer / developer & Writes and maintains software \\
      IT support / technician & Installs, maintains, and troubleshoots \\
      \bottomrule
    \end{tabular}
  \end{center}
  \vspace{0.2cm}
  \muted{All of these people together are the humanware of the system.}
\end{frame}

\begin{frame}{Users, Operators, and Administrators}
  \begin{itemize}
    \item \key{User (end user)}: the person who uses the finished system to
      perform a task --- a clerk, a student, a customer.
    \item \key{Operator}: runs the computer system, feeds and monitors jobs,
      and handles routine operation.
    \item \key{System administrator}: manages servers, user accounts, and
      access rights.
    \item \key{Network / database administrator}: manages the network or the
      database.
    \item \muted{Administrators manage the system; operators run it; users
      use its output.}
  \end{itemize}
\end{frame}

\begin{frame}{Programmers and IT Support}
  \begin{itemize}
    \item \key{Programmer / developer}: writes, tests, and maintains the
      software the system runs.
    \item \key{IT support / help desk}: answers user problems and
      troubleshoots faults.
    \item \key{Technician}: installs, repairs, and maintains the hardware.
    \item \muted{All of these are humanware --- the human element, not
      machinery.}
  \end{itemize}
  \begin{block}{One idea}
    Humanware is about people and their roles, not about devices or programs.
  \end{block}
\end{frame}

\begin{frame}{Hardware + Software + Humanware}
  \begin{center}
    \begin{tabular}{p{0.24\textwidth}p{0.64\textwidth}}
      \toprule
      \textbf{Component} & \textbf{What it is} \\
      \midrule
      Hardware & The physical parts of the computer \\
      Software & The programs and instructions \\
      Humanware & The people who use and manage them \\
      \bottomrule
    \end{tabular}
  \end{center}
  \vspace{0.3cm}
  \begin{block}{All three are required}
    Hardware without software does nothing, and both are useless without the
    people who run and use them.
  \end{block}
\end{frame}

\begin{frame}{IT in Governance: The Role}
  \begin{itemize}
    \item \key{Governance} is the process of decision-making and its
      implementation in an organisation or a state.
    \item \key{ICT} (Information and Communication Technology) lets government
      deliver services over computers and networks.
    \item The \key{role of IT in governance} is to make government work
      faster, cheaper, more transparent, and more accessible.
    \item \muted{This is the ``role of IT in governance'' annexure topic.}
  \end{itemize}
\end{frame}

\begin{frame}{Benefits of IT in Governance}
  \begin{itemize}
    \item \key{Transparency}: citizens can see rules, status, and decisions.
    \item \key{Accountability}: actions are recorded and traceable.
    \item \key{Efficiency}: less paperwork and faster processing.
    \item \key{Convenience}: services available from home, round the clock.
    \item Lower cost of service delivery and reduced opportunities for
      corruption.
  \end{itemize}
  \begin{alertblock}{Takeaway}
    The aim is \key{citizen-centric} service --- the citizen at the centre of
    every process.
  \end{alertblock}
\end{frame}

\begin{frame}{IT Governance vs E-Governance (Near-Neighbour)}
  \begin{center}
    \begin{tabular}{p{0.26\textwidth}p{0.62\textwidth}}
      \toprule
      \textbf{Term} & \textbf{Meaning} \\
      \midrule
      IT governance & Directing and controlling an organisation's own IT
        resources --- governance \emph{of} IT \\
      E-governance & Using IT to govern and deliver public services --- IT
        \emph{in} governance \\
      \bottomrule
    \end{tabular}
  \end{center}
  \vspace{0.3cm}
  \begin{alertblock}{Trap alert}
    IT governance is \key{governance of IT}; e-governance is \key{IT in
    governance}. Do not swap the two.
  \end{alertblock}
\end{frame}

\begin{frame}{E-Governance Basics: What It Is}
  \begin{itemize}
    \item \key{E-governance} is the application of ICT by government to
      deliver services, share information, and transact with citizens,
      businesses, and other government agencies.
    \item It uses the internet, databases, and digital records instead of
      paper and counters.
    \item It makes government services faster, wider-reaching, and easier to
      audit.
    \item \muted{Named platforms and G2x matching are covered in detail in
      Lesson 57.}
  \end{itemize}
  \begin{block}{Definition}
    E-governance is the use of information and communication technology to
    improve the way government serves and interacts with the public.
  \end{block}
\end{frame}

\begin{frame}{G2C, G2G, G2B, G2E}
  \begin{center}
    \begin{tabular}{p{0.16\textwidth}p{0.30\textwidth}p{0.44\textwidth}}
      \toprule
      \textbf{Type} & \textbf{Full form} & \textbf{Interaction} \\
      \midrule
      G2C & Government to Citizen & Services to the general public \\
      G2G & Government to Government & Between government departments \\
      G2B & Government to Business & With companies and traders \\
      G2E & Government to Employee & With government staff \\
      \bottomrule
    \end{tabular}
  \end{center}
  \vspace{0.25cm}
  \muted{These four are the standard interaction types of e-governance.}
\end{frame}

\begin{frame}{Digital Public Services}
  \begin{itemize}
    \item Online delivery of government services: certificates, licences,
      land records, bill payment, applications, and grievance redressal.
    \item Citizens reach them through service portals and Common Service
      Centres (CSCs).
    \item Typical features: online application, status tracking, digital
      payment, and downloadable documents.
    \item \muted{A service is a ``digital'' public service when the citizen
      can complete the task without visiting an office in person.}
  \end{itemize}
  \begin{block}{One idea}
    Digital public services move a government task from a paper counter to an
    online process.
  \end{block}
\end{frame}

\begin{frame}{Electronic Records}
  \begin{itemize}
    \item An \key{electronic record} is information created, stored, and
      retrieved in electronic (digital) form.
    \item Examples: e-documents, e-files, scanned papers, database entries,
      and email.
    \item Benefits: paperless storage, easy search and retrieval, sharing,
      and backup.
    \item \muted{An electronic record can have legal validity when it is
      authenticated, for example with a digital signature.}
  \end{itemize}
\end{frame}

\begin{frame}{Authentication: Verifying Identity}
  \begin{itemize}
    \item \key{Authentication} is verifying that a user, device, or system is
      who or what it claims to be.
    \item It answers the question: \key{``Who are you?''}
    \item Common methods: password, PIN, OTP, biometrics, and a digital
      certificate.
    \item \muted{Authentication must succeed before the system grants
      access.}
  \end{itemize}
  \begin{block}{Definition}
    Authentication is the process of confirming a claimed identity before
    access is allowed.
  \end{block}
\end{frame}

\begin{frame}{Authentication vs Authorization (Near-Neighbour)}
  \begin{center}
    \begin{tabular}{p{0.24\textwidth}p{0.64\textwidth}}
      \toprule
      \textbf{Point} & \textbf{Meaning} \\
      \midrule
      Authentication & Proves identity --- ``who are you?'' \\
      Authorization & Grants or denies permission --- ``what may you do?'' \\
      \bottomrule
    \end{tabular}
  \end{center}
  \vspace{0.3cm}
  \begin{alertblock}{Trap alert}
    Authentication \key{identifies}; authorization \key{permits}.
    Authentication comes first, then authorization.
  \end{alertblock}
\end{frame}

\begin{frame}{Authentication Factors (2FA and MFA)}
  \begin{itemize}
    \item \key{Something you know}: a password or PIN.
    \item \key{Something you have}: an OTP on your phone, a token, or a smart
      card.
    \item \key{Something you are}: a fingerprint, iris, or face.
    \item \key{Two-factor authentication (2FA)} uses two \emph{different}
      factors; \key{multi-factor authentication (MFA)} uses two or more.
    \item \muted{Two passwords are not two factors --- they are the same
      factor twice.}
  \end{itemize}
\end{frame}

\begin{frame}{Trap Alert: Facts to Keep Straight}
  \begin{itemize}
    \item Humanware = \key{people}; it is neither hardware nor software.
    \item IT governance = governance \key{of} IT; e-governance = IT \key{in}
      governance.
    \item G2C = Government to Citizen (not ``Government to Company'').
    \item Authentication = identity; authorization = permission.
    \item An electronic record is digital information, not a printed paper
      file.
  \end{itemize}
\end{frame}

\begin{frame}{Lesson 47: Recall}
  \begin{itemize}
    \item \key{Humanware} (liveware / peopleware) is the people element:
      users, operators, administrators, programmers, and IT support.
    \item Hardware, software, and humanware are all required for a working
      system.
    \item IT in governance improves transparency, accountability,
      efficiency, and citizen convenience.
    \item IT governance = governance of IT; e-governance = IT in governance;
      G2x = G2C, G2G, G2B, G2E.
    \item Digital public services, electronic records, and authentication are
      the building blocks of e-governance.
    \item Authentication proves identity; authorization grants permission.
  \end{itemize}
\end{frame}
\end{document}
